\properlane{protected-pilgrimage}{God guards his people on an unfinished journey}

The Collect asks freedom from all adversity. The Offertory answers with
Job's losses and a grievous wound, and Communion still cries for help.
Those words would contradict one another if protection meant that no
hostile action ever reached God's people. They form a coherent prayer
when God's guardianship is understood through the journey of salvation:
real deliverance, a still-dangerous passage, sustaining gifts, and a
final rest beyond the enemy's reach. The prayer for immediate rescue
belongs fully within that journey. It is not a mistake to be corrected
by telling the sufferer to want pain instead.

\subsection{The permission does not make the malice good}
Job's opening chapters allow the reader to know something Job himself
does not see: the accuser has challenged the sincerity of his fidelity.
Would Job fear God without the shelter around his prosperity? The first
permission concerns what he possesses; the next permits bodily attack
but protects his life. The Offertory condenses those scenes. Its
uncompromising opening---simple, upright, fearing God---prevents a
reading in which the subsequent wound proves that the first description
was false.

Gregory begins with the adversary's request. Satan tells God to put
forth his hand because he cannot strike by autonomous power. Even
hostility to the creator does not make a creature self-subsisting.
Gregory distinguishes the enemy's evil will from the permitted power
through which it acts: the desire to injure is Satan's, while the
limits of the action remain under divine government. Permission is
not approval of the creature's intention.\footnote{Gregory,
\work{Moralia} II.10.16--18.}

That distinction changes the meaning of the Introit's sovereignty.
Mordecai's ``Lord of all'' is not a distant ruler who has ceded a
province of creation to an equally powerful evil lord. Nor does the
claim make Haman's murderous project righteous. An endangered people
can ask the creator to frustrate a wicked purpose precisely because
the purpose does not possess final dominion. The prayer is confident
without being an assertion that no threat is present.

Gregory gives close attention to the limits of the permission. God
lets the enemy loose and reins him in. In the first trial Job's person
is withheld; in the second his life remains protected. Gregory
describes the same man as being in the hand of the adversary and in
the hand of his helper, but in different respects. The wound does
not prove the helper's hand has vanished.\footnote{Gregory,
\work{Moralia} II.11.19; III.5.6--7.}

The distinction is more exact than saying that suffering is only an
appearance. Job's loss is dreadful and bodily. Gregory's consolation
is that evil cannot determine the whole person and the final outcome
simply by doing bodily harm. His exposition supplies a way to resist
despair; it supplies no warrant for refusing help to a person whose
danger can be relieved. The petition for protection and the work of
protecting another belong on the same side.

\subsection{What ``without cause'' rules out}
The sharper difficulty is God's statement that Job has been afflicted
``without cause.'' Gregory refuses two easy answers. One would make
God speak falsely by discovering a secret punitive guilt in Job. The
other would make God act without justice or purpose. He distinguishes
the absence of a crime requiring punishment from the presence of a
providential end in which hidden virtue becomes manifest and patience
grows. Job is afflicted without the first kind of cause; his trial
is not without the second.\footnote{Gregory, \work{Moralia} III.3.3.}

Gregory also examines the phrase in which Satan is said to have
``moved'' God against Job. God has not become enraged under the devil's
persuasion. Human speech describes divine action in the terms through
which human beings recognize an action: because we often strike when
moved, the strike is called a movement. The unchanging God does not
pass from wisdom into emotional darkness. The account preserves
God's justice and truth together rather than imagining him defeated
in an argument with his creature.\footnote{Gregory,
\work{Moralia} III.3.4.}

The Gradual gives this answer a different expression. Before the
mountains, through the generations, God \emph{is}. Augustine attends
to that present tense and identifies the refuge as eternal being
rather than one temporal power among others. People begin to take
refuge in God; he does not begin to be God when they arrive. Their
changing circumstances therefore cannot be read directly as the
history of a deity growing strong and weak.\footnote{Augustine,
Ps.~89 \S\S2--4.}

These claims do not reveal the purpose of every particular loss.
Gregory has an inspired narrative which expressly calls Job upright
and discloses the heavenly challenge. A neighbor confronted with
bereavement or illness does not thereby possess that disclosure about
the sufferer. The text rules out automatic diagnosis of guilt and
teaches trust in divine government. It does not authorize announcing
that this person's wound was necessary for improvement.

The Sunday Secret places restoration in God's powerful goodness.
Human endurance alone is not the saving remedy. The offering seeks
propitiation and salvation from the One who has willed them. Schuster's
typological reading of Job as Christ brings the innocent sufferer
toward that sacrifice: the sinless Christ voluntarily bears our sins.
His reading belongs to fulfillment in Christ; it does not vindicate
the friends' literal accusation against Job.\footnote{Schuster,
\work{Sacramentary} III, pp.~178,181.}

\subsection{Why people who have escaped Egypt still need armor}
The Alleluia announces departure: Israel came out of Egypt. In the
psalm's continuation, Judah becomes God's sanctuary and Israel his
dominion. The rescued people belongs to the deliverer. Augustine
follows Paul's reading of the sea as a baptismal figure and the
wilderness food and drink as spiritual gifts. Yet Paul tells the
same baptized Church that many who received those signs fell in the
wilderness. The sea is a genuine passage out of bondage, and it is
not the final rest.\footnote{Augustine, Ps.~113 \S\S3--6,
following 1 Cor.~10:1--6.}

Chrysostom reaches the same journey by a different path: the girdle
of truth in Ephesians. Girding makes a traveler ready to move. Israel
ate Passover with loins girded, shoes on, and staff at hand. The
historical posture reminded its children that their deliverance was
about to begin. Chrysostom then identifies Christ as the Christian
Passover and asks his hearers to eat as people ready for their own
departure. The image turns Eucharistic reception outward into a
manner of living.\footnote{Chrysostom, \work{Ephesians} XXIII,
on 6:14 and the Passover exposition.}

He does not leave the comparison at a single resemblance. The
Christian community has crossed the sea and remains in a wilderness;
the grace of the Spirit supplies light and shelter; its nourishment
is greater than manna; Christ leads its Exodus. The destination is
heavenly rest, where warfare ends. The gifts arrive before that rest,
because the travelers need them on the way. Chrysostom particularly
opposes ingratitude, indolence, and attachment to wealth: a traveler
who treats the camp as his permanent city has misunderstood both
the food and the journey.

This sequence explains the armor's place without making the Gospel
of peace absurd. Resistance is needed because the pilgrim has been
freed for God and can still be tempted away. Chrysostom's breastplate
is a virtuous life; the faith which meets the fiery darts rests on
the promised good. The evil day has to be faced, not renamed a good
day. But its threats are not the measure of the destination.\footnote{Chrysostom, \work{Ephesians} XXII--XXIV, especially XXIII's
contrast between wilderness warfare and heavenly rest.}

The Gospel supplies a severe check upon a merely comforting use of
this pilgrimage. The servant has truly been released, and still
falls. His danger comes not only from an external enemy but from his
own deliberate refusal of mercy. Perseverance includes charity; it
cannot be reduced to surviving hostility while becoming hostile
oneself. The first reading followed the causal claim of the king's
gift upon conduct. Here the same event marks a stage of salvation:
departure is not arrival, and the delivered person must continue in
the life received.

\subsection{The Church can be helped without every martyr escaping death}
The Communion's cry is stronger than a declaration of general
confidence. It asks when God will execute judgment and begs him to
help against unjust persecution. Augustine hears the whole Christ,
head and members, in the psalm's longing. Saints before the
Incarnation desired Christ's coming; saints after the Ascension
desire his manifestation as judge. The many lives are joined by
one hope in the promised Savior.\footnote{Augustine, Ps.~118,
Caph, on vv.~81--82.}

His reading of fainting is particularly important. The full biblical
verse says the soul has fainted after salvation. Augustine distinguishes
the direction of that failing: it is longing \emph{for} the good,
not departure \emph{from} it. Thus emotional exhaustion is not
automatically faithlessness. Hope rests on God's word while the
desired consolation is delayed. The chant itself omits the verb
of fainting, but retains the salvation, hope, and petition which
his exposition places within that larger context.

Augustine next hears the martyrs' cry from the Apocalypse in the
question about judgment. The Church is not promised a quiet
interval in which she disappears from history and waits elsewhere
for the end. Her testimony continues amid persecution until the
judgment she requests. Bellarmine likewise refuses to turn the
question into calculation of a date. It expresses longing for
the security of everlasting peace.\footnote{Augustine, Ps.~118
Caph, on v.~84; Bellarmine, Ps.~118:81--88, especially 84.}

The strongest test arrives when Augustine asks whether the prayer
for help was in vain if martyrs were killed. He follows the psalm's
next words: the people did not forsake God's commandments, and
asked to be quickened by his mercy. Their life was preserved at
the level at which denial of truth would have lost it. They lived
through dying for the truth. This is not a denial of death. It is
an account of a salvation that death cannot overthrow.\footnote{Augustine, Ps.~118, Caph, on vv.~86--88.}

The Sunday Collect's request can therefore be prayed without
pretending that every bodily danger has already been removed.
It seeks real deliverance, and the final completeness of that
deliverance exceeds present bodily security. The Sunday
Postcommunion gives the destination its name: immortality.
The food is given now, and its inward fruit is sought now,
while its promised horizon extends beyond the present conflict.

Schuster adds a distinct present emphasis to the Communion.
He reads rejection of light and truth as a judgment already
borne by their enemies. That does not answer the entire
``when?'' in Augustine or Bellarmine, whose expectation includes
final judgment and rest. Present privation and future
consummation are related, but they are different claims.
The assembly's petition can carry both the seriousness of
present refusal and the still-awaited completion of justice.\footnote{Schuster, \work{Sacramentary} III, p.~181.}

\subsection{The journey is common, sacramental, and open to all peoples}
The Collect's family does not travel as a collection of sealed
individuals. The Trinity Preface joins its present thanksgiving
to angelic praise and confesses the God who is its destination.
The Persons' distinction does not divide divine saving power.
Augustine's account of the Son and Spirit being sent without
inferiority illuminates that confession: the gifts bestowed in
time belong to the equal divine majesty which the prayer adores.

When the mission prayers are added, travelers also become
workers. The Collect asks that the word run, that it be made
known, and that all peoples know God and the sent Christ.
The journey is neither a retreat from other people nor a
campaign to treat them as spiritual enemies. Its courage is
the confidence to speak the truth; its end is communion in
the God it proclaims.

The mission Secret stretches worship from sunrise to sunset.
Christ's self-giving for all makes that universal horizon
possible. The mission Postcommunion, however, does not treat
reception of redemption as a finished possession requiring
nothing further. It asks continuing progress in true faith
through the gift of everlasting salvation. The three prayers
give sending, offering, and growth their place within the
pilgrimage of a community already nourished.

Luke's commemoration, where admitted, supplies a concrete
apostolic pattern. The body bears the cross for God's name;
the heavenly gifts heal and make service free; the received
altar gift sanctifies the soul and safeguards it. The Collect
and Postcommunion must be heard together. A saint known in
prayer as a cross-bearer does not exemplify a safety defined
as absence of every wound. The Secret's glory names a further
horizon for its cure, while 2 Corinthians~4 describes the
outward affliction and inward renewal of apostolic ministry.

This reading leaves the sufferer's plea intact. Mordecai
asks rescue; Job's wounds are grievous; the psalmist asks
help. Faith is not a measurable covering whose failure
would prove that a victim believed too little. It receives
God's gifts, resists the temptation to abandon truth and
charity, and hopes for the rest in which persecution truly
ends. The Mass supplies armor and nourishment because
the road is still real.

\begin{fourSenses}
\item[Literal.] Threatened Israel petitions its creator; Paul commands
steadfastness; upright Job suffers within limits on the adversary;
the psalmist remains persecuted; and a pardoned servant can still
refuse the life his release requires.
\item[Allegorical.] The Church makes her Exodus in Christ. Augustine's
baptismal passage and Chrysostom's Passover, Spirit, and food identify
a people delivered and still journeying; Schuster reads Job's innocent
suffering toward the crucified Christ.
\item[Moral.] Ask and work for rescue, stand in truth and peace, refuse
to diagnose another's wounds as proof of guilt, persevere in mercy,
and let the sacramental gift shape the whole interior life.
\item[Anagogical.] Christ's body awaits judgment and unassailable rest.
Immortal nourishment and thanksgiving with the angels orient the
pilgrimage to enduring communion with the triune God.
\end{fourSenses}

\clearpage
